\documentclass[11pt]{article}

\usepackage[margin=1in]{geometry}
\usepackage[hidelinks]{hyperref}
\usepackage{xcolor}
\usepackage{enumitem}
\usepackage{parskip}
\usepackage{array}

\definecolor{primarycolor}{RGB}{0, 102, 204}

\hypersetup{colorlinks=true, urlcolor=primarycolor, linkcolor=primarycolor}

\setlist[itemize]{leftmargin=1.5em, itemsep=2pt, topsep=2pt}

\newcommand{\code}[1]{\texttt{\small #1}}

% Colored, bold section titles without titlesec.
\makeatletter
\renewcommand{\section}{\@startsection{section}{1}{\z@}%
  {-3.0ex \@plus -1ex \@minus -.2ex}{1.5ex \@plus.2ex}%
  {\Large\bfseries\color{primarycolor}}}
\renewcommand{\subsection}{\@startsection{subsection}{2}{\z@}%
  {-2.5ex\@plus -1ex \@minus -.2ex}{1.0ex \@plus .2ex}%
  {\large\bfseries}}
\makeatother

\begin{document}

%---------- TITLE ----------
\begin{center}
    {\Huge \bfseries \color{primarycolor} Review Intelligence} \\[4pt]
    {\large An End-to-End NLP Project: from pretrained transformers to a fine-tuned, deployed model} \\[8pt]
    {\large Rahiq Raees} \\[2pt]
    {\normalsize July 2026} \\[6pt]
    \normalsize
    \href{https://github.com/rahiqraees/review-intelligence}{GitHub Repository} $\quad|\quad$
    \href{https://huggingface.co/spaces/rahiqr/review-intelligence}{Live Demo} $\quad|\quad$
    \href{https://huggingface.co/rahiqr/distilbert-amazon-sentiment}{Fine-tuned Model}
\end{center}

\vspace{6pt}
\hrule
\vspace{10pt}

\tableofcontents
\vspace{10pt}
\hrule
\vspace{10pt}

%---------- EXECUTIVE SUMMARY ----------
\section{Executive Summary}

\textbf{Review Intelligence} is an end-to-end natural language processing (NLP)
system that turns unstructured customer reviews into structured, actionable
signal: \textbf{sentiment}, \textbf{topics}, and a short \textbf{summary}.

The project was built deliberately as one evolving system rather than a set of
disconnected tutorials. It progresses from off-the-shelf transformer models,
through a deep understanding of how those models work internally, to a
\textbf{custom model fine-tuned on real product-review data}, and finally to a
\textbf{live, publicly hosted web application} that serves that model.

Key outcomes:
\begin{itemize}
    \item Fine-tuned a DistilBERT transformer that improved product-review
    sentiment accuracy from \textbf{88.6\% to 89.7\%} (F1 0.899) over a strong
    pretrained baseline, with an honest, controlled evaluation.
    \item Published the fine-tuned model to the Hugging Face Hub and deployed an
    interactive demo on Hugging Face Spaces that serves it live.
    \item Engineered the codebase to a production standard: an installable
    Python package, an automated test suite, and continuous integration (CI).
\end{itemize}

%---------- MOTIVATION ----------
\section{Motivation and Objective}

Businesses receive customer feedback as free-form text at a volume no team can
read manually. Automatically classifying each review's sentiment, identifying
which topics it touches (shipping, product quality, price, support), and
summarizing it are the building blocks of a feedback dashboard that tells a
company \emph{what} customers are unhappy about and \emph{how much}.

The objective was to build such a system while demonstrating the full modern NLP
workflow --- using pretrained transformer models, understanding their internals,
adapting them to a specific domain through fine-tuning, and deploying the result
--- all with the software-engineering rigor expected in production
(reproducible environments, testing, CI, and deployment).

%---------- WHAT IT DOES ----------
\section{What the System Does}

The system exposes three analyzers, each deliberately built on a different
transformer \emph{architecture}. This mapping is a core NLP concept: different
model designs suit different tasks.

\begin{center}
\renewcommand{\arraystretch}{1.4}
\begin{tabular}{|>{\raggedright\arraybackslash}p{3.2cm}|>{\raggedright\arraybackslash}p{4.6cm}|>{\raggedright\arraybackslash}p{5.2cm}|}
\hline
\textbf{Capability} & \textbf{Task} & \textbf{Transformer architecture} \\
\hline
Sentiment scoring & Text classification & Encoder (DistilBERT) --- reads the whole
sentence to \emph{understand} it \\
\hline
Topic detection & Zero-shot classification & Encoder (BART fine-tuned on natural
language inference) \\
\hline
Summarization & Sequence-to-sequence & Encoder--decoder (DistilBART) --- reads
the review and \emph{generates} a shorter one \\
\hline
\end{tabular}
\end{center}

\textbf{Example.} Given the review \emph{``The product itself is fine, but it took
three weeks to arrive and the box was crushed,''} the system returns sentiment
\code{NEGATIVE}, top topic \emph{shipping and delivery}, and a one-line summary.

%---------- TECHNICAL APPROACH ----------
\section{Technical Approach}

The system was built in four stages, each adding a capability and a deeper level
of understanding.

\subsection{Stage 1 --- Baseline with pretrained pipelines}

The first version used Hugging Face \code{pipeline()} objects --- a high-level
interface to run production-grade transformer models with a single call. This
established a working baseline on real reviews and demonstrated the three
transformer families (encoder, decoder, encoder--decoder).

A notable technique here is \textbf{zero-shot classification} for topic
detection: rather than training a classifier on labeled topic data, the system
reframes ``is this review about shipping?'' as a natural-language-inference
question. This lets the topic labels be defined at run time and changed freely
without any retraining.

\subsection{Stage 2 --- Understanding the internals}

To move beyond treating the models as black boxes, the sentiment pipeline was
reproduced manually, stage by stage:
\begin{itemize}
    \item \textbf{Tokenizer:} converts text into integer token IDs and adds the
    special tokens the model expects, producing an ``attention mask'' that marks
    real tokens versus padding.
    \item \textbf{Model:} runs the transformer to produce raw scores (logits).
    \item \textbf{Post-processing:} applies a softmax to convert logits into
    probabilities, then selects the highest-scoring label.
\end{itemize}
The manual result matched the pipeline's output exactly, confirming a correct
understanding of every step --- a prerequisite for the fine-tuning that follows.

\subsection{Stage 3 --- Fine-tuning a domain-specific model}

The baseline sentiment model is trained on \emph{movie} reviews (the SST-2
dataset), while this project analyzes \emph{product} reviews --- a different
domain with different vocabulary. To close this gap, a DistilBERT model was
\textbf{fine-tuned} (an application of transfer learning) on a streamed,
class-balanced subset of the Amazon product-reviews dataset.

Methodology:
\begin{itemize}
    \item Data loaded and processed with the Hugging Face \code{datasets}
    library; sequences tokenized in batches with dynamic padding for efficiency.
    \item Training run with the Hugging Face \code{Trainer} API (6{,}000
    training reviews, 2 epochs, learning rate $2\times10^{-5}$).
    \item Evaluated honestly: the fine-tuned model and the movie-trained baseline
    were both scored on the \emph{same} 1{,}000 held-out product reviews, using
    accuracy and F1.
\end{itemize}

\subsection{Stage 4 --- Publishing and deployment}

The fine-tuned model was published to the Hugging Face Hub as a model repository
with a model card documenting its data, results, and usage. An interactive web
application was built with Gradio and deployed to Hugging Face Spaces, where it
runs on hosted CPU hardware and serves the fine-tuned model to any user with the
link.

%---------- RESULTS ----------
\section{Results}

\begin{center}
\renewcommand{\arraystretch}{1.4}
\begin{tabular}{|>{\raggedright\arraybackslash}p{7cm}|c|c|}
\hline
\textbf{Model (evaluated on product reviews)} & \textbf{Accuracy} & \textbf{F1} \\
\hline
Baseline (movie-trained, SST-2) & 88.6\% & 0.888 \\
\hline
\textbf{Fine-tuned (product reviews)} & \textbf{89.7\%} & \textbf{0.899} \\
\hline
\end{tabular}
\end{center}

\textbf{Interpretation.} The improvement of $+1.1$ percentage points is real but
modest, and understanding \emph{why} is part of the result. Binary sentiment
transfers well across domains --- ``I love it'' reads the same in a movie or a
product review --- so the off-the-shelf baseline was already strong, leaving a
small domain gap to close. Reporting a genuine, well-measured improvement against
a strong baseline, and understanding its size, reflects sound evaluation
practice: a controlled comparison on a held-out set, rather than an inflated
number from a weak baseline or a leaky test.

%---------- ENGINEERING ----------
\section{Software Engineering Practices}

The project was engineered like production software, not a one-off notebook:
\begin{itemize}
    \item \textbf{Installable package} with a clean, reusable API separating
    library code from runnable scripts.
    \item \textbf{Reproducible environment} pinned via a Conda environment file,
    including a deliberate version pin to keep results stable across library
    updates.
    \item \textbf{Automated tests} (\code{pytest}), split into fast unit tests
    (no model downloads, safe for CI) and slower integration tests.
    \item \textbf{Continuous integration} via GitHub Actions, running the test
    suite across Python 3.10 and 3.11 on every change.
    \item \textbf{Deployment} of both a model artifact (to the Hub) and a live
    application (to Spaces).
\end{itemize}

%---------- SKILLS ----------
\section{Skills and Tools Demonstrated}

\begin{itemize}
    \item \textbf{NLP / Transformers:} tokenization, model internals, encoder vs.
    encoder--decoder architectures, zero-shot classification, fine-tuning,
    transfer learning.
    \item \textbf{Ecosystem:} Hugging Face Transformers, Datasets, Hub, and
    Spaces; PyTorch; Gradio.
    \item \textbf{Evaluation:} controlled baseline comparison, accuracy and F1,
    train/test hygiene.
    \item \textbf{Engineering:} Python packaging, \code{pytest}, GitHub Actions
    CI, reproducible environments, model and application deployment.
\end{itemize}

%---------- LIMITATIONS ----------
\section{Limitations and Future Work}

\begin{itemize}
    \item \textbf{Binary sentiment.} The current model predicts positive or
    negative only; genuinely mixed reviews are forced into one class. A natural
    next step is a three-class (negative / neutral / positive) model to capture
    nuance.
    \item \textbf{Demo hardware.} The live demo runs on free CPU hardware, so the
    first request is slow while models load, then fast.
    \item \textbf{Next project.} The same skills are being extended to a
    retrieval-augmented question-answering (RAG) system in a new domain,
    combining semantic search with a question-answering model.
\end{itemize}

%---------- LINKS ----------
\section{Links and Reproducibility}

\begin{itemize}
    \item \textbf{Source code:}
    \href{https://github.com/rahiqraees/review-intelligence}{github.com/rahiqraees/review-intelligence}
    \item \textbf{Live demo:}
    \href{https://huggingface.co/spaces/rahiqr/review-intelligence}{huggingface.co/spaces/rahiqr/review-intelligence}
    \item \textbf{Fine-tuned model:}
    \href{https://huggingface.co/rahiqr/distilbert-amazon-sentiment}{huggingface.co/rahiqr/distilbert-amazon-sentiment}
\end{itemize}

The repository includes setup instructions, runnable scripts for each stage, and
the recorded evaluation metrics, so the full pipeline can be reproduced from
scratch.

\end{document}
